\chapter{Concepts I --- Constraints, requires-Clauses, and requires-Expressions}
\label{ch:c20-concepts-i}

\begin{quote}
\emph{Concepts are the first of C++20's four pillars and the one that most directly changes day-to-day template code. They turn template parameters from unchecked ``duck typing resolved deep inside instantiation'' into named, compile-time-verified contracts. This chapter covers the mechanics: defining a concept, attaching constraints with a requires-clause, expressing ad-hoc requirements with a requires-expression, and the four ways to spell a constrained template.}
\end{quote}

Before C++20, expressing ``this template only works for types that support \texttt{+}'' meant SFINAE, \texttt{std::enable\_if}, tag dispatch, or a buried \texttt{static\_assert} --- all producing diagnostics measured in screenfuls. Concepts replace that machinery with a first-class language feature that names the requirement, checks it at the call site, and emits a one-line error when it fails.

\section{The Problem Concepts Solve}
\label{sec:c20-concepts-problem}

Hand a \texttt{std::list} to a random-access algorithm and pre-C++20 the compiler complains \textbf{deep inside the algorithm}, dozens of instantiation frames from the call site. The \textbf{bouncer} analogy: old C++ admits everyone and discovers the problem inside; concepts check IDs at the door.

\begin{itemize}
\item \textbf{Diagnostics} become short and point at the call site.
\item \textbf{Overload resolution} can select by type properties cleanly.
\item \textbf{Interfaces} become self-documenting --- the constraint \emph{is} the contract.
\end{itemize}

There is \textbf{no runtime cost}: concepts are evaluated at compile time and erased.

\section{Defining a Concept}
\label{sec:c20-defining-concept}

A \textbf{concept} is a named compile-time predicate over template parameters, introduced with the \texttt{concept} keyword.

\begin{lstlisting}[language=C++,caption={Defining concepts from a trait and from a requires-expression}]
#include <concepts>
#include <type_traits>

// Form 1: built on a type trait
template<typename T>
concept Integral = std::is_integral_v<T>;

// Form 2: built on a requires-expression (tests expression validity)
template<typename T>
concept Addable = requires(T a, T b) {
    { a + b } -> std::convertible_to<T>;   // a + b must be valid AND convertible to T
};
\end{lstlisting}

Concepts compose with \texttt{\&\&} and \texttt{||}, which (unlike a plain \texttt{bool}) participate in \textbf{subsumption} (Chapter~33).

\begin{lstlisting}[language=C++,caption={Composing concepts}]
template<typename T>
concept Number = std::integral<T> || std::floating_point<T>;

template<typename T>
concept SignedNumber = Number<T> && std::is_signed_v<T>;
\end{lstlisting}

\section{The Four Ways to Constrain a Template}
\label{sec:c20-four-syntaxes}

\begin{lstlisting}[language=C++,caption={The four constraint syntaxes, all equivalent in effect}]
#include <concepts>

// (1) Trailing requires-clause
template<typename T>
T inc1(T x) requires std::integral<T> { return x + 1; }

// (2) Leading requires-clause
template<typename T>
    requires std::integral<T>
T inc2(T x) { return x + 1; }

// (3) Constrained template parameter (concept in place of 'typename')
template<std::integral T>
T inc3(T x) { return x + 1; }

// (4) Abbreviated function template (constrained auto parameter)
std::integral auto inc4(std::integral auto x) { return x + 1; }
\end{lstlisting}

\begin{center}
\begin{tabular}{ll}
\hline
\textbf{Syntax} & \textbf{Best when} \\
\hline
Trailing \texttt{requires} & constraint depends on parameter types or is complex \\
Leading \texttt{requires} & standalone clause you want prominent \\
\texttt{template<C T>} & single simple concept on one parameter \\
\texttt{C auto} & terse generic functions \\
\hline
\end{tabular}
\end{center}

\section{requires-Clauses}
\label{sec:c20-requires-clause}

A \textbf{requires-clause} attaches a boolean constraint expression to a template. Its operand must be a concept, a parenthesized \texttt{bool} constant expression, or a conjunction/disjunction of these.

\begin{lstlisting}[language=C++,caption={requires-clauses with composed and parenthesized constraints}]
#include <concepts>
#include <type_traits>

template<typename T>
    requires std::integral<T> && (sizeof(T) >= 4)   // concept && parenthesized bool
T checked_double(T x) { return x * 2; }

template<typename T>
    requires std::copyable<T>
class Box { T value; };
\end{lstlisting}

A \textbf{raw boolean expression} like \texttt{sizeof(T) >= 4} must be \textbf{parenthesized} when combined with \texttt{\&\&}/\texttt{||}, because only primary expressions are allowed there.

\section{requires-Expressions}
\label{sec:c20-requires-expression}

A \textbf{requires-expression} yields a \texttt{bool} at compile time: \texttt{true} if every requirement is satisfied. Its optional parameter list introduces fictional variables, never evaluated.

\begin{lstlisting}[language=C++,caption={Anatomy of a requires-expression}]
template<typename T>
concept Stack = requires(T s, typename T::value_type v) {
    s.push(v);          // simple requirement
    s.pop();
    { s.top() } -> std::same_as<typename T::value_type&>;  // compound requirement
    typename T::value_type;                                 // type requirement
    requires std::default_initializable<T>;                 // nested requirement
};
\end{lstlisting}

\subsection{Simple Requirements}
\label{sec:c20-simple-req}

A \textbf{simple requirement} is an expression followed by \texttt{;}, satisfied if the expression is \emph{valid} (it is never evaluated).

\begin{lstlisting}[language=C++,caption={Simple requirements}]
template<typename T>
concept Incrementable = requires(T x) {
    x++;        // postfix increment must be valid
    ++x;        // prefix increment must be valid
    x + x;      // addition must be valid
};
\end{lstlisting}

\subsection{Type Requirements}
\label{sec:c20-type-req}

A \textbf{type requirement} is \texttt{typename} followed by a type name, satisfied if the named type exists.

\begin{lstlisting}[language=C++,caption={Type requirements}]
template<typename T>
concept HasValueType = requires {
    typename T::value_type;
    typename T::iterator;
};
\end{lstlisting}

\subsection{Compound Requirements}
\label{sec:c20-compound-req}

A \textbf{compound requirement} wraps an expression in braces and optionally asserts \texttt{noexcept} and/or a return-type constraint.

\begin{lstlisting}[language=C++,caption={Compound requirements with return-type constraints and noexcept}]
#include <concepts>

template<typename T>
concept Hashable = requires(const T& x) {
    { std::hash<T>{}(x) } -> std::convertible_to<std::size_t>;
};

template<typename T>
concept NothrowSwappable = requires(T& a, T& b) {
    { a.swap(b) } noexcept;
};
\end{lstlisting}

\texttt{\{ expr \} -> Concept;} means ``\texttt{expr} is valid, and \texttt{Concept<decltype((expr))>} holds.''

\subsection{Nested Requirements}
\label{sec:c20-nested-req}

A \textbf{nested requirement} is \texttt{requires} followed by a constraint expression; its boolean value must be \texttt{true}.

\begin{lstlisting}[language=C++,caption={Nested requirement enforces a predicate, not just validity}]
#include <concepts>

template<typename T>
concept SmallTrivial = requires {
    requires std::is_trivially_copyable_v<T>;   // the predicate must be TRUE
    requires sizeof(T) <= 16;                    // and this one too
};
\end{lstlisting}

\texttt{sizeof(T) <= 16;} as a \emph{simple} requirement only checks validity (useless); \texttt{requires sizeof(T) <= 16;} checks truth.

\section{Constrained auto}
\label{sec:c20-constrained-auto}

A concept can replace \texttt{auto} in parameters, return types, and variables.

\begin{lstlisting}[language=C++,caption={Constrained auto in parameters, returns, and variables}]
#include <concepts>
#include <ranges>

void take_integer(std::integral auto x);          // abbreviated function template
std::integral auto make_index() { return 0u; }    // constrained return type
std::ranges::range auto first_n(auto&& r);         // constrained return
std::integral auto n = compute();                  // constrained variable
\end{lstlisting}

Each \texttt{Concept auto} parameter introduces an \textbf{independent} template parameter; for identical types use a named parameter.

\section{Constraint-Based Overloading vs. SFINAE}
\label{sec:c20-overloading}

\begin{lstlisting}[language=C++,caption={Overloading on concepts replaces enable\_if}]
#include <concepts>
#include <iostream>

void describe(std::integral auto x)       { std::cout << "integer: " << x << '\n'; }
void describe(std::floating_point auto x) { std::cout << "float: "   << x << '\n'; }

int main() {
    describe(42);     // integer
    describe(3.14);   // float
}
\end{lstlisting}

When multiple constrained overloads are viable, the \textbf{most constrained} wins via subsumption (Chapter~33).

\section{Diagnostics: The Real Payoff}
\label{sec:c20-diagnostics}

\begin{lstlisting}[language=C++,caption={A constraint violation produces a localized, readable error}]
#include <concepts>
#include <list>
#include <ranges>

void needs_random_access(std::ranges::random_access_range auto&& r);

int main() {
    std::list<int> lst{1, 2, 3};
    needs_random_access(lst);
    // error: constraints not satisfied
    //   note: 'random_access_range<std::list<int>&>' was not satisfied
    //   note: 'std::list<int>::iterator' does not satisfy 'random_access_iterator'
}
\end{lstlisting}

This is the difference between a two-minute fix and a two-hour spelunk.

\section{Professional Insights}
\label{sec:c20-concepts-i-insights}

\textbf{Constrain every public template interface.} The cost is one concept name; the benefit is a call-site diagnostic and a self-documenting contract. An unconstrained public template parameter is now a code smell.

\textbf{Prefer named concepts to inline requires-expressions in interfaces.} Named concepts are reusable, appear in diagnostics by name, and participate in subsumption so overloads order correctly.

\textbf{Know the simple-vs-nested requirement trap cold.} \texttt{sizeof(T) <= 16;} checks validity, not truth --- almost always a silent bug. The predicate form is \texttt{requires sizeof(T) <= 16;}.

\textbf{Use compound requirements to constrain results, not just calls.} \texttt{\{ expr \} -> std::convertible\_to<U>;} pins down what an expression yields, far stronger than the bare \texttt{expr;}.

\textbf{Remember there is no runtime cost and no ABI footprint.} Concepts generate no code and add no indirection --- precisely why they belong in hot-path template libraries.
